\documentclass[a4paper,10pt]{article}
\usepackage[utf8]{inputenc}
\usepackage[T1]{fontenc}
\usepackage[french]{babel}
\usepackage{amsmath,amssymb}
\usepackage{geometry}
\usepackage{array,tabularx}
\usepackage{tikz}
\geometry{margin=1.6cm}
\pagestyle{empty}
\newcommand{\ligne}{\par\vspace{0.55cm}\noindent\dotfill\par}
\newcommand{\carreaux}[1]{\begin{center}\begin{tikzpicture}\draw[step=0.5cm,gray!35,very thin] (0,0) grid (17,#1);\end{tikzpicture}\end{center}}
\newenvironment{exercice}[2]{\paragraph{Exercice #1 \hfill #2 points}\mbox{}\par}{\medskip}
\begin{document}
\begin{center}\Large\bfseries Devoir surveillé 19 - Algorithme et programmation\end{center}
\vspace{2mm}
\begin{exercice}{1}{4}Écrire un algorithme qui demande un nombre et affiche son triple augmenté de $7$.\end{exercice}\begin{exercice}{2}{5}Un programme effectue : choisir un nombre, ajouter $4$, multiplier par $2$. Tester avec $3$ puis écrire l'expression en fonction de $x$.\end{exercice}\begin{exercice}{3}{5}Dans Scratch, un lutin répète $4$ fois : avancer de $80$, tourner de $90^\circ$. Quelle figure trace-t-il ?\end{exercice}\begin{exercice}{4}{6}On veut tracer un triangle équilatéral. Proposer les blocs de déplacement et de rotation à répéter.\end{exercice}\end{document}
